%%%%%%%%%%%%%%%%%%%%%%%%%%%%%%%%%%%%%%%%%%%%%%%%%%%%%%%%%%%%%%%%%%%%%%%%%%%%%%%%
%   Copyright 2009, Oracle and/or its affiliates.
%   All rights reserved.
%
%
%   Use is subject to license terms.
%
%   This distribution may include materials developed by third parties.
%
%%%%%%%%%%%%%%%%%%%%%%%%%%%%%%%%%%%%%%%%%%%%%%%%%%%%%%%%%%%%%%%%%%%%%%%%%%%%%%%%

\chapter{Numbers}
\chaplabel{lib:numbers-basic}


\section{Rational Numbers}
\seclabel{basic-rationals}

The trait \EXP{\mathbb{Q}} (\STR{QQ}) encompasses all finite rational numbers,
the results of dividing any integer by any nonzero integer,
and three more elements:
\EXP{+\infty} and \EXP{-\infty} (written \EXP{1/0} and \EXP{-1/0}),
the results of dividing a positive or a negative integer by zero,
and the indefinite rational (written \EXP{0/0}), which is used as the result
of dividing zero by zero or of adding \EXP{-\infty} to \EXP{+\infty}.

The number types of this chapter are
the integer types \EXP{\mathbb{Z}}, \EXP{\mathbb{Z}64}, \EXP{\mathbb{Z}32},
\EXP{\mathbb{N}64}, and \EXP{\mathbb{N}32} (\secref{basic-integers}),
the rational type \EXP{\mathbb{Q}},
and the floating-point types \EXP{\mathbb{R}64} (\STR{RR64}) and \EXP{\mathbb{R}32} (\STR{RR32}).
They are siblings under the trait \TYP{Number}:
each is a subtype of \TYP{Number}, and none is a subtype of another.
Some conversions between them are coercions (\chapref{conversions-coercions}),
each declared by the type converted to;
since coercions do not chain (\secref{coercion}), each pair is declared.
Those among the integer types are given in \secref{basic-integers};
the others are these:

\begin{tabbing}
\EXP{\mathbb{Q}} (\STR{QQ}) coerces from \EXP{\mathbb{Z}}, \EXP{\mathbb{Z}64}, \EXP{\mathbb{Z}32}, \EXP{\mathbb{N}64}, and \EXP{\mathbb{N}32}.
\\
\EXP{\mathbb{R}64} (\STR{RR64}) coerces from \EXP{\mathbb{R}32}, from \EXP{\mathbb{Z}32}, and from integer numerals (\secref{literals}).
\end{tabbing}

Every other conversion between number types is written explicitly;
in particular, a conversion into \EXP{\mathbb{R}64} from
\EXP{\mathbb{Z}}, \EXP{\mathbb{Z}64}, \EXP{\mathbb{N}64}, \EXP{\mathbb{N}32}, or \EXP{\mathbb{Q}}
is explicit, written \EXP{\VAR{asFloat}(x)} with the method that \library\ declare on \TYP{Number}.
\revision{revival-rr32}{Before this revision the passage named \EXP{\mathbb{R}64}
as the one floating-point type, and \library\ declared \EXP{\mathbb{R}32} a
subtype of it, so that a value of \EXP{\mathbb{R}32} was also a value of
\EXP{\mathbb{R}64}, which \secref{types-libraries} excludes.
It now names \EXP{\mathbb{R}32} beside \EXP{\mathbb{R}64}, a sibling under
\TYP{Number} as \library\ now declare it, with the coercion into
\EXP{\mathbb{R}64} that the example of \secref{widening} relies on.}

It is often desirable to know that a value lies in a subset of its type,
such as the positive values, the nonzero values, or the finite values.
Such subsets are not types of their own: a program tests a value at run time,
with the comparison operators for its sign and, for a floating-point value,
with the two checks that \library\ declare on \EXP{\mathbb{R}64}:
%getter check(): Maybe[\RR64\]
%getter check_star(): Maybe[\RR64\]
\begin{Fortress}
\(\KWD{getter}\:\VAR{check}()\COLON \TYP{Maybe}\llbracket\mathbb{R}64\rrbracket\)\\
\(\KWD{getter}\:\VAR{check}^*()\COLON \TYP{Maybe}\llbracket\mathbb{R}64\rrbracket\)
\end{Fortress}
The getter \VAR{check} returns \EXP{\TYP{Just}(x)} for a finite value \VAR{x}
and \TYP{Nothing} otherwise;
the getter \EXP{\VAR{check}^*} returns \EXP{\TYP{Just}(x)} for a value \VAR{x}
that is not a NaN and \TYP{Nothing} otherwise.

Here we present the trait \EXP{\mathbb{Q}} and its methods.
Because \EXP{\mathbb{Q}} holds \EXP{0/0}, which is unordered with respect to every
rational value, \EXP{\mathbb{Q}} is neither totally ordered nor a field.
The original design of the rational types, in which \EXP{\mathbb{Q}} and
seventeen refined rational types are instantiations of one trait, is kept as a
superseded design in \secref{advanced-rationals}.
\revision{revival-numbers}{The Working Draft of February 2011 divided the
rationals into eighteen types: \EXP{\mathbb{Q}}, the type \EXP{\mathbb{Q}^*} that
adds the two infinities, the type \EXP{\mathbb{Q}^{\#}} that adds \EXP{0/0}, and
five refinements of each by sign, each a subtype of others and of the
corresponding real types, and it said that the type system tracks them through
arithmetic.
\Library\ never declared them, and the trait by which
\secref{advanced-rationals} defines them is refused by instantiation exclusion
(\secref{revival-rule}).
The text above describes the number types the library declares, siblings under
\TYP{Number}, with the coercions it declares; a subset such as the positive
values is tested at run time.
\EXP{\mathbb{Q}} holds the values that \EXP{\mathbb{Q}^*} and \EXP{\mathbb{Q}^{\#}}
added, as the library's rational type does and as \chapref{operators} already
says (``For rational results, division by zero produces \EXP{1/0}'').
Because it holds \EXP{0/0}, the listing below no longer gives it the field and
order traits of the original design, which that design itself withheld from
\EXP{\mathbb{Q}^{\#}}.}

%% trait QQ
%%     extends { Number, StandardPartialOrder[\QQ\], StandardMinMax[\QQ\],
%%               AdditiveGroup[\QQ\], MultiplicativeRing[\QQ\] }
%%   coerce(x: ZZ)
%%   coerce(x: ZZ64)
%%   coerce(x: ZZ32)
%%   coerce(x: NN64)
%%   coerce(x: NN32)
%%   coerce(x: Identity[\+\])
%%   coerce(x: Identity[\DOT\])
%%   coerce(x: Identity[\TIMES\])
%%   coerce(x: Identity[\juxtaposition\])
%%   coerce(x: Zero[\DOT\])
%%   coerce(x: Zero[\TIMES\])
%%   coerce(x: Zero[\juxtaposition\])
%%   opr juxtaposition(self, other: QQ): QQ
%%   opr +(self): QQ
%%   opr +(self, other: QQ): QQ
%%   opr -(self): QQ
%%   opr -(self, other: QQ): QQ
%%   opr DOT(self, other: QQ): QQ
%%   opr TIMES(self, other: QQ): QQ
%%   opr /(self): QQ
%%   opr /(self, other: QQ): QQ
%%   opr ^(self, power: ZZ): QQ
%%   opr <(self, other: QQ): Boolean
%%   opr <=(self, other: QQ): Boolean
%%   opr =(self, other: QQ): Boolean
%%   opr >=(self, other: QQ): Boolean
%%   opr >(self, other: QQ): Boolean
%%   opr CMP(self, other: QQ): Comparison
%%   opr MAX(self, other: QQ): QQ
%%   opr MIN(self, other: QQ): QQ
%%   opr MAXNUM(self, other: QQ): QQ
%%   opr MINNUM(self, other: QQ): QQ
%%   opr |self| : QQ
%%   signum(self): ZZ
%%   numerator(self): ZZ
%%   denominator(self): ZZ
%%   floor(self): ZZ
%%   opr |\ self /|: ZZ
%%   ceiling(self): ZZ
%%   opr |/ self \|: ZZ
%%   round(self): ZZ
%%   truncate(self): ZZ
%%   opr ||\ self /||: ZZ
%%   opr ||/ self \||: ZZ
%%   opr |||\ self /|||: ZZ
%%   opr |||/ self \|||: ZZ
%%   realpart(self): QQ
%%   imagpart(self): QQ
%% end
\begin{Fortress}
\(\KWD{trait}\:\mathbb{Q}\)\\
{\tt~~~~}\pushtabs\=\+\(    \KWD{extends} \{\, \null\)\pushtabs\=\+\(\TYP{Number}, \TYP{StandardPartialOrder}\llbracket\mathbb{Q}\rrbracket, \TYP{StandardMinMax}\llbracket\mathbb{Q}\rrbracket,\)\\
\(              \TYP{AdditiveGroup}\llbracket\mathbb{Q}\rrbracket, \TYP{MultiplicativeRing}\llbracket\mathbb{Q}\rrbracket\,\}\)\-\-\\\poptabs\poptabs
{\tt~~}\pushtabs\=\+\(  \KWD{coerce}(x\COLON \mathbb{Z})\)\\
\(  \KWD{coerce}(x\COLON \mathbb{Z}64)\)\\
\(  \KWD{coerce}(x\COLON \mathbb{Z}32)\)\\
\(  \KWD{coerce}(x\COLON \mathbb{N}64)\)\\
\(  \KWD{coerce}(x\COLON \mathbb{N}32)\)\\
\(  \KWD{coerce}(x\COLON \TYP{Identity}\llbracket+\rrbracket)\)\\
\(  \KWD{coerce}(x\COLON \TYP{Identity}\llbracket\cdot\rrbracket)\)\\
\(  \KWD{coerce}(x\COLON \TYP{Identity}\llbracket\times\rrbracket)\)\\
\(  \KWD{coerce}(x\COLON \TYP{Identity}\llbracket\KWD{juxtaposition}\rrbracket)\)\\
\(  \KWD{coerce}(x\COLON \TYP{Zero}\llbracket\cdot\rrbracket)\)\\
\(  \KWD{coerce}(x\COLON \TYP{Zero}\llbracket\times\rrbracket)\)\\
\(  \KWD{coerce}(x\COLON \TYP{Zero}\llbracket\KWD{juxtaposition}\rrbracket)\)\\
\(  \KWD{opr}\;\;\KWD{juxtaposition}(\KWD{self}, \VAR{other}\COLON \mathbb{Q})\COLON \mathbb{Q}\)\\
\(  \KWD{opr} +(\KWD{self})\COLON \mathbb{Q}\)\\
\(  \KWD{opr} +(\KWD{self}, \VAR{other}\COLON \mathbb{Q})\COLON \mathbb{Q}\)\\
\(  \KWD{opr} -(\KWD{self})\COLON \mathbb{Q}\)\\
\(  \KWD{opr} -(\KWD{self}, \VAR{other}\COLON \mathbb{Q})\COLON \mathbb{Q}\)\\
\(  \KWD{opr} \mathord{\cdot}(\KWD{self}, \VAR{other}\COLON \mathbb{Q})\COLON \mathbb{Q}\)\\
\(  \KWD{opr} \mathord{\times}(\KWD{self}, \VAR{other}\COLON \mathbb{Q})\COLON \mathbb{Q}\)\\
\(  \KWD{opr} /(\KWD{self})\COLON \mathbb{Q}\)\\
\(  \KWD{opr} /(\KWD{self}, \VAR{other}\COLON \mathbb{Q})\COLON \mathbb{Q}\)\\
\(  \KWD{opr} \mathord{\hbox{\tt\char'136}}(\KWD{self}, \VAR{power}\COLON \mathbb{Z})\COLON \mathbb{Q}\)\\
\(  \KWD{opr} \mathord{<}(\KWD{self}, \VAR{other}\COLON \mathbb{Q})\COLON \TYP{Boolean}\)\\
\(  \KWD{opr} \mathord{\leq}(\KWD{self}, \VAR{other}\COLON \mathbb{Q})\COLON \TYP{Boolean}\)\\
\(  \KWD{opr} \mathord{=}(\KWD{self}, \VAR{other}\COLON \mathbb{Q})\COLON \TYP{Boolean}\)\\
\(  \KWD{opr} \mathord{\geq}(\KWD{self}, \VAR{other}\COLON \mathbb{Q})\COLON \TYP{Boolean}\)\\
\(  \KWD{opr} \mathord{>}(\KWD{self}, \VAR{other}\COLON \mathbb{Q})\COLON \TYP{Boolean}\)\\
\(  \KWD{opr} \mathord{\OPR{CMP}}(\KWD{self}, \VAR{other}\COLON \mathbb{Q})\COLON \TYP{Comparison}\)\\
\(  \KWD{opr} \mathord{\OPR{MAX}}(\KWD{self}, \VAR{other}\COLON \mathbb{Q})\COLON \mathbb{Q}\)\\
\(  \KWD{opr} \mathord{\OPR{MIN}}(\KWD{self}, \VAR{other}\COLON \mathbb{Q})\COLON \mathbb{Q}\)\\
\(  \KWD{opr} \mathord{\OPR{MAXNUM}}(\KWD{self}, \VAR{other}\COLON \mathbb{Q})\COLON \mathbb{Q}\)\\
\(  \KWD{opr} \mathord{\OPR{MINNUM}}(\KWD{self}, \VAR{other}\COLON \mathbb{Q})\COLON \mathbb{Q}\)\\
\(  \KWD{opr} \left|\mathord{\KWD{self}}\right| \mathrel{\mathtt{:}}\:\mathbb{Q}\)\\
\(  \VAR{signum}(\KWD{self})\COLON \mathbb{Z}\)\\
\(  \VAR{numerator}(\KWD{self})\COLON \mathbb{Z}\)\\
\(  \VAR{denominator}(\KWD{self})\COLON \mathbb{Z}\)\\
\(  \VAR{floor}(\KWD{self})\COLON \mathbb{Z}\)\\
\(  \KWD{opr} \lfloor \mathord{\KWD{self}} \rfloor\COLON \mathbb{Z}\)\\
\(  \VAR{ceiling}(\KWD{self})\COLON \mathbb{Z}\)\\
\(  \KWD{opr} \lceil \mathord{\KWD{self}} \rceil\COLON \mathbb{Z}\)\\
\(  \VAR{round}(\KWD{self})\COLON \mathbb{Z}\)\\
\(  \VAR{truncate}(\KWD{self})\COLON \mathbb{Z}\)\\
\(  \KWD{opr} \lhfloor \mathord{\KWD{self}} \rhfloor\COLON \mathbb{Z}\)\\
\(  \KWD{opr} \lhceil \mathord{\KWD{self}} \rhceil\COLON \mathbb{Z}\)\\
\(  \KWD{opr} \lhhfloor \mathord{\KWD{self}} \rhhfloor\COLON \mathbb{Z}\)\\
\(  \KWD{opr} \lhhceil \mathord{\KWD{self}} \rhhceil\COLON \mathbb{Z}\)\\
\(  \VAR{realpart}(\KWD{self})\COLON \mathbb{Q}\)\\
\(  \VAR{imagpart}(\KWD{self})\COLON \mathbb{Q}\)\-\\\poptabs
\(\KWD{end}\)
\end{Fortress}
\revision{revival-rationals}{The Working Draft of February 2011 declared
\EXP{\mathbb{Q}} a subtype of the real numbers and of \EXP{\mathbb{Q}^*}, a
field with the nonzero rationals \EXP{\mathbb{Q}_{\neq}} as its multiplicative
group, and totally ordered; it typed the results of the reciprocal, of division
and of exponentiation as \EXP{\mathbb{Q}^*} or \EXP{\mathbb{Q}^{\#}}, those of
the absolute value and of the four hyperfloor and hyperceiling operators as
\EXP{\mathbb{Q}_{\geq}} and \EXP{\mathbb{N}}, and it declared seventeen check
methods, each returning a refined type or throwing \TYP{CastError}.
The listing and the method entries below take the type \library\ declare where
they declare the method, and \EXP{\mathbb{Q}} or \EXP{\mathbb{Z}} where they do
not; the two comparisons by \OPR{CMP} become the one the library declares.
The check methods leave the listing: the library declares \VAR{check} and
\EXP{\VAR{check}^*} on \EXP{\mathbb{R}64} (\secref{basic-rationals}) and no
check by sign.
What the method entries said ``for types \EXP{\mathbb{Q}^*} and
\EXP{\mathbb{Q}^{\#}}'' they now say of \EXP{\mathbb{Q}}, which holds those
values.
The coercions from the five integer types are new, with the method entry that
follows.}
The four traits that the listing's \KWD{extends} clause names after \TYP{Number}
are those \library\ declare on \EXP{\mathbb{Q}}; their laws hold for the finite
rationals and need not hold where \EXP{0/0} or an infinity takes part, as the laws
of the same traits on \EXP{\mathbb{R}64} need not hold where a NaN takes part.
\revision{revival-rational-algebra}{Before this revision the listing extended
\TYP{Number} alone, the field and order traits of the original design having left
it (\secref{revival-rationals}); it now names the traits \library\ declare on
\EXP{\mathbb{Q}}, and the sentence above says where their laws hold.}


%%   coerce(x: ZZ)
%%   coerce(x: ZZ64)
%%   coerce(x: ZZ32)
%%   coerce(x: NN64)
%%   coerce(x: NN32)
\Method{\EXP{\KWD{coerce}(\VAR{x}\COLON \mathbb{Z})}}
\Method*{\EXP{\KWD{coerce}(\VAR{x}\COLON \mathbb{Z}64)}}
\Method*{\EXP{\KWD{coerce}(\VAR{x}\COLON \mathbb{Z}32)}}
\Method*{\EXP{\KWD{coerce}(\VAR{x}\COLON \mathbb{N}64)}}
\Method*{\EXP{\KWD{coerce}(\VAR{x}\COLON \mathbb{N}32)}}

Each of these coercions converts an integer to the rational number with the same value.


%%   coerce(x: Identity[\+\])
%%   coerce(x: Identity[\DOT\])
%%   coerce(x: Identity[\TIMES\])
%%   coerce(x: Identity[\juxtaposition\])
%%   coerce(x: Zero[\DOT\])
%%   coerce(x: Zero[\TIMES\])
%%   coerce(x: Zero[\juxtaposition\])
\Method{\EXP{\KWD{coerce}(\VAR{x}\COLON \TYP{Identity}\llbracket+\rrbracket)}}
\Method*{\EXP{\KWD{coerce}(\VAR{x}\COLON \TYP{Identity}\llbracket\cdot\rrbracket)}}
\Method*{\EXP{\KWD{coerce}(\VAR{x}\COLON \TYP{Identity}\llbracket\times\rrbracket)}}
\Method*{\EXP{\KWD{coerce}(\VAR{x}\COLON \TYP{Identity}\llbracket\KWD{juxtaposition}\rrbracket)}}
\Method*{\EXP{\KWD{coerce}(\VAR{x}\COLON \TYP{Zero}\llbracket\cdot\rrbracket)}}
\Method*{\EXP{\KWD{coerce}(\VAR{x}\COLON \TYP{Zero}\llbracket\times\rrbracket)}}
\Method*{\EXP{\KWD{coerce}(\VAR{x}\COLON \TYP{Zero}\llbracket\KWD{juxtaposition}\rrbracket)}}

The identity for \EXP{+} is 0.

The identity for \KWD{juxtaposition} or \EXP{\cdot} or \EXP{\times} is 1.

The zero for \KWD{juxtaposition} or \EXP{\cdot} or \EXP{\times} is 0.



%  opr juxtaposition(self, other: QQ): QQ
\Method{\EXP{\KWD{opr}\ \KWD{juxtaposition}(\KWD{self}, \VAR{other}\COLON \mathbb{Q})\COLON \mathbb{Q}}}

Juxtaposition of rational expressions is equivalent to using the multiplication operator \EXP{\cdot}.


%  opr +(self): QQ
\Method{\EXP{\KWD{opr} +(\KWD{self})\COLON \mathbb{Q}}}

The unary addition operator \EXP{+} simply returns its argument.


%  opr +(self, other: QQ): QQ
\Method{\EXP{\KWD{opr} +(\KWD{self}, \VAR{other}\COLON \mathbb{Q})\COLON \mathbb{Q}}}

The binary addition operator \EXP{+} returns the sum of its arguments.

The sum of an infinity
and either a finite rational or another infinity of the same sign is equal to the given infinity,
but the sum of infinities of differing sign is \EXP{0/0},
and the sum of \EXP{0/0} and any rational value is \EXP{0/0}.


%  opr -(self): QQ
\Method{\EXP{\KWD{opr} -(\KWD{self})\COLON \mathbb{Q}}}

The unary negation operator \EXP{-} returns the negative of its argument.

The negative of
\EXP{+\infty} is \EXP{-\infty}, the negative of \EXP{-\infty} is \EXP{+\infty},
and the negative of \EXP{0/0} is \EXP{0/0}.


%  opr -(self, other: QQ): QQ
\Method{\EXP{\KWD{opr} -(\KWD{self}, \VAR{other}\COLON \mathbb{Q})\COLON \mathbb{Q}}}

The binary subtraction operator \EXP{-} returns the difference of its arguments,
which is equal to the sum of (a) the first argument and (b) the negation of the second argument.


%%  opr DOT(self, other: QQ): QQ
%%  opr TIMES(self, other: QQ): QQ
\Method{\EXP{\KWD{opr} \mathord{\cdot}(\KWD{self}, \VAR{other}\COLON \mathbb{Q})\COLON \mathbb{Q}}}
\Method*{\EXP{\KWD{opr} \mathord{\times}(\KWD{self}, \VAR{other}\COLON \mathbb{Q})\COLON \mathbb{Q}}}

The multiplication operator \EXP{\cdot} (\STR{DOT}) returns the product of its arguments.
The multiplication operator \EXP{\times} (\STR{TIMES}) does exactly the same thing.

The product of
\EXP{0/0} and any rational value is \EXP{0/0}, and the product of
zero and an infinity (regardless of sign) is \EXP{0/0};
the product of an infinity and any rational value other than zero and \EXP{0/0}
is an infinity whose sign is positive if and only if the two arguments have the same sign.


%  opr /(self): QQ
\Method{\EXP{\KWD{opr} /(\KWD{self})\COLON \mathbb{Q}}}

The unary reciprocal operator \EXP{/} returns the reciprocal of its argument.
The reciprocal of zero is \EXP{+\infty}.

The reciprocal of
either \EXP{+\infty} or \EXP{-\infty} is zero, and the reciprocal of \EXP{0/0} is \EXP{0/0}.


%  opr /(self, other: QQ): QQ
\Method{\EXP{\KWD{opr} /(\KWD{self}, \VAR{other}\COLON \mathbb{Q})\COLON \mathbb{Q}}}

The binary division operator \EXP{/} returns the quotient of its arguments,
which is equal to the product of (a) the first argument and (b) the reciprocal of the second argument.


%  opr ^(self, power: ZZ): QQ
\Method{\EXP{\KWD{opr} \mathord{\hbox{\tt\char'136}}(\KWD{self}, \VAR{power}\COLON \mathbb{Z})\COLON \mathbb{Q}}}

Exponentiation of a rational number to an integer power produces a rational result.
If the \VAR{power} is \EXP{0}, then the result is always \EXP{1}, even if the rational number base is \EXP{0}
(this definition is somewhat arbitrary but is computationally useful).

%%  property FORALL(x, y:ZZ) x^y = 1/(x^(-y))
%%  property FORALL(x, y:ZZ) x^y = x^(|\ y/2 /|) x^(|/ y/2 \|)
\begin{Fortress}
{\tt~}\pushtabs\=\+\( \KWD{property} \forall(x, y\COLONOP\mathbb{Z})\; x^{y} = 1/(x^{-y})\)\\
\( \KWD{property} \forall(x, y\COLONOP\mathbb{Z})\; x^{y} = x^{(\lfloor y/2 \rfloor)} x^{(\lceil y/2 \rceil)}\)\-\\\poptabs
\end{Fortress}


%  opr <(self, other: QQ): Boolean
%  opr <=(self, other: QQ): Boolean
%  opr =(self, other: QQ): Boolean
%  opr >=(self, other: QQ): Boolean
%  opr >(self, other: QQ): Boolean
\Method{\EXP{\KWD{opr} \mathord{<}(\KWD{self}, \VAR{other}\COLON \mathbb{Q})\COLON \TYP{Boolean}}}
\Method*{\EXP{\KWD{opr} \mathord{\leq}(\KWD{self}, \VAR{other}\COLON \mathbb{Q})\COLON \TYP{Boolean}}}
\Method*{\EXP{\KWD{opr} \mathord{=}(\KWD{self}, \VAR{other}\COLON \mathbb{Q})\COLON \TYP{Boolean}}}
\Method*{\EXP{\KWD{opr} \mathord{\geq}(\KWD{self}, \VAR{other}\COLON \mathbb{Q})\COLON \TYP{Boolean}}}
\Method*{\EXP{\KWD{opr} \mathord{>}(\KWD{self}, \VAR{other}\COLON \mathbb{Q})\COLON \TYP{Boolean}}}

The comparison operators \EXP{<}, \EXP{\leq}, \EXP{=}, \EXP{\geq}, and \EXP{>}
allow any rational value to be compared numerically to any other rational value.

The rational values are totally ordered except for \EXP{0/0},
which is unordered with respect to all other rational values;
moreover, for compatibility with floating-point arithmetic, \EXP{0/0} is
unordered with respect to itself, and therefore these five comparison operators
always return \VAR{false} if either argument is \EXP{0/0}.
The value \EXP{-\infty} is less than any finite rational value,
and \EXP{+\infty} is greater than any finite rational value.

For \EXP{\neq} \see{operator-NE}.

%  opr CMP(self, other: QQ): Comparison
\Method{\EXP{\KWD{opr} \mathord{\OPR{CMP}}(\KWD{self}, \VAR{other}\COLON \mathbb{Q})\COLON \TYP{Comparison}}}

The \OPR{CMP} operator compares the arguments and returns one of the four values
\TYP{LessThan}, \TYP{EqualTo}, \TYP{GreaterThan}, and \TYP{Unordered}.


%  opr MAX(self, other: QQ): QQ
%  opr MIN(self, other: QQ): QQ
%  opr MAXNUM(self, other: QQ): QQ
%  opr MINNUM(self, other: QQ): QQ
\Method{\EXP{\KWD{opr} \mathord{\OPR{MAX}}(\KWD{self}, \VAR{other}\COLON \mathbb{Q})\COLON \mathbb{Q}}}
\Method*{\EXP{\KWD{opr} \mathord{\OPR{MIN}}(\KWD{self}, \VAR{other}\COLON \mathbb{Q})\COLON \mathbb{Q}}}
\Method*{\EXP{\KWD{opr} \mathord{\OPR{MAXNUM}}(\KWD{self}, \VAR{other}\COLON \mathbb{Q})\COLON \mathbb{Q}}}
\Method*{\EXP{\KWD{opr} \mathord{\OPR{MINNUM}}(\KWD{self}, \VAR{other}\COLON \mathbb{Q})\COLON \mathbb{Q}}}

The operators \OPR{MAX} and \OPR{MAXNUM} return whichever argument is larger
in the total order defined by \EXP{<}, \EXP{\leq}, \EXP{=}, \EXP{\geq}, \EXP{>}, and \OPR{CMP},
and the operators \OPR{MIN} and \OPR{MINNUM} return whichever argument is smaller.
(For all four, if the arguments are equal, then the result equals that same value.)

The operators \OPR{MAXNUM} and \OPR{MINNUM} differ from
\OPR{MAX} and \OPR{MIN} in their treatment of \EXP{0/0}:
if one argument is \EXP{0/0} and the other is not, then \OPR{MAX} or \OPR{MIN} returns \EXP{0/0} but
\OPR{MAXNUM} or \OPR{MINNUM} returns the argument that is not \EXP{0/0}.


%  opr |self| : QQ
\Method{\EXP{\KWD{opr} \left|\mathord{\KWD{self}}\right| \mathrel{\mathtt{:}} \mathbb{Q}}}

The absolute value operator \EXP{\left|\ldots\right|} returns the negative of this rational number
if the argument is less than zero, and otherwise returns the argument.

The absolute value of \EXP{0/0} is \EXP{0/0}.


%  signum(self): ZZ
\Method{\EXP{\VAR{signum}(\KWD{self})\COLON \mathbb{Z}}}

The method \VAR{signum} returns \EXP{-1} if this rational number is less than zero,
\EXP{0} if this rational number is zero, and \EXP{1} if this rational number is greater than zero.

The signum of \EXP{0/0} is \EXP{0/0}.


%  numerator(self): ZZ
%  denominator(self): ZZ
\Method{\EXP{\VAR{numerator}(\KWD{self})\COLON \mathbb{Z}}}
\Method*{\EXP{\VAR{denominator}(\KWD{self})\COLON \mathbb{Z}}}

The method \VAR{numerator} returns the numerator of this rational number,
and the method \VAR{denominator} returns the denominator of this rational number,
when this rational number is represented in lowest terms (such that the greatest common divisor
of numerator and denominator is 1).

The numerator of \EXP{+\infty} is \EXP{1},
the numerator of \EXP{-\infty} is \EXP{-1},
and the numerator of \EXP{0/0} is \EXP{0};
the denominator of \EXP{+\infty}, \EXP{-\infty}, or \EXP{0/0} is \EXP{0}.


%  floor(self): ZZ
%  opr LF self RF: ZZ
%  ceiling(self): ZZ
%  opr LC self RC: ZZ
%  truncate(self): ZZ
%  round(self): ZZ
\Method{\EXP{\VAR{floor}(\KWD{self})\COLON \mathbb{Z}}}
\Method*{\EXP{\KWD{opr} \lfloor \mathord{\KWD{self}} \rfloor\COLON \mathbb{Z}}}
\Method*{\EXP{\VAR{ceiling}(\KWD{self})\COLON \mathbb{Z}}}
\Method*{\EXP{\KWD{opr} \lceil \mathord{\KWD{self}} \rceil\COLON \mathbb{Z}}}
\Method*{\EXP{\VAR{round}(\KWD{self})\COLON \mathbb{Z}}}
\Method*{\EXP{\VAR{truncate}(\KWD{self})\COLON \mathbb{Z}}}

The method \VAR{floor}, likewise the enclosing operator \EXP{\lfloor\ldots\rfloor},
returns the largest integer that is not greater than this rational number.

The method \VAR{ceiling}, likewise the enclosing operator \EXP{\lceil\ldots\rceil},
returns the smallest integer that is not less than this rational number.

The method \VAR{round} returns the integer that is closest to this rational number,
but if this rational number is exactly halfway between two consecutive integers,
then \VAR{round} returns whichever of the two integers is even.

The method \VAR{truncate} returns the ceiling of this rational number if it is
negative, and otherwise returns the floor of this rational number.
(This has the effect of taking the floor of the magnitude, also called
``rounding toward zero.'')

All of these methods throw a \TYP{DivisionByZero} if this rational number is
\EXP{+\infty}, \EXP{-\infty}, or \EXP{0/0}, none of which has an integer value;
so does a division by zero whose result is an integer (\chapref{operators}).
\revision{revival-rational-rounding}{Before this revision the sentence said that
all of these methods simply return the argument if it is \EXP{+\infty},
\EXP{-\infty}, or \EXP{0/0}, a rational number where the entries above declare
an integer result. It now says that they throw a \TYP{DivisionByZero} there, as
\library\ do.}


%%   opr ||\ self /||: ZZ
%%   opr ||/ self \||: ZZ
%%   opr |||\ self /|||: ZZ
%%   opr |||/ self \|||: ZZ
\begin{Fortress}
{\tt~~}\pushtabs\=\+\(  \KWD{opr} \lhfloor \mathord{\KWD{self}} \rhfloor\COLON \mathbb{Z}\)\\
\(  \KWD{opr} \lhceil \mathord{\KWD{self}} \rhceil\COLON \mathbb{Z}\)\\
\(  \KWD{opr} \lhhfloor \mathord{\KWD{self}} \rhhfloor\COLON \mathbb{Z}\)\\
\(  \KWD{opr} \lhhceil \mathord{\KWD{self}} \rhhceil\COLON \mathbb{Z}\)\-\\\poptabs
\end{Fortress}

The hyperfloor operation \EXP{\lhfloor{}x\rhfloor} computes \EXP{2^{(\lfloor {log}_{2} x \rfloor)}}
and returns the result as a natural number.
If the argument is equal to \EXP{0}, the result is \EXP{0}.
If the argument is negative, an \TYP{InvalidArgument} is thrown.

The hyperceiling operation \EXP{\lhceil{}x\rhceil} computes \EXP{2^{(\lceil {log}_{2} x \rceil)}}
and returns the result as a natural number.
If the argument is equal to \EXP{0}, the result is \EXP{0}.
If the argument is negative, an \TYP{InvalidArgument} is thrown.

The hyperhyperfloor operation \EXP{\lhhfloor{}x\rhhfloor} computes \EXP{2^{(\lhfloor {log}_{2} x \rhfloor)}}
and returns the result as a natural number.
If the argument is equal to \EXP{0} or \EXP{1}, the result is the same as the argument.
If the argument is negative, an \TYP{InvalidArgument} is thrown.

The hyperhyperceiling operation \EXP{\lhhceil{}x\rhhceil} computes \EXP{2^{(\lhceil {log}_{2} x \rhceil)}}
and returns the result as a natural number.
If the argument is equal to \EXP{0} or \EXP{1}, the result is the same as the argument.
If the argument is negative, an \TYP{InvalidArgument} is thrown.


%  realpart(self): QQ
\Method{\EXP{\VAR{realpart}(\KWD{self})\COLON \mathbb{Q}}}

The method \VAR{realpart} for a rational number simply returns its argument.


%  imagpart(self): QQ
\Method{\EXP{\VAR{imagpart}(\KWD{self})\COLON \mathbb{Q}}}

The method \VAR{imagpart} for a rational number simply returns zero.

